% =====================================================================
%  preamble.tex -- packages, layout, notation and planning helpers
% =====================================================================

% ---------- layout ----------------------------------------------------
\usepackage[a4paper,margin=25mm,bindingoffset=0mm]{geometry}
\usepackage[T1]{fontenc}
\usepackage{lmodern}
\usepackage{microtype}
\usepackage{setspace}
\onehalfspacing
\usepackage[english]{babel}
\usepackage{csquotes}

% ---------- maths -----------------------------------------------------
\usepackage{amsmath,amssymb,amsthm,mathtools,bm}
\numberwithin{equation}{chapter}
\theoremstyle{plain}
\newtheorem{proposition}{Proposition}[chapter]
\theoremstyle{definition}
\newtheorem{definition}{Definition}[chapter]
\theoremstyle{remark}
\newtheorem{remark}{Remark}[chapter]

% ---------- tables, figures, units -------------------------------------
\usepackage{graphicx}
\graphicspath{{figures/}}
\usepackage{booktabs,tabularx,multirow,array,makecell}
\usepackage{siunitx}
\sisetup{group-separator={,},group-minimum-digits=4}
\usepackage{caption,subcaption}
\captionsetup{font=small,labelfont=bf}
\usepackage{xcolor}
\usepackage{tikz}
\usetikzlibrary{arrows.meta,positioning,fit,calc,shapes.geometric}
\usepackage{enumitem}
\setlist{itemsep=2pt,topsep=4pt}

% ---------- references --------------------------------------------------
% IEEE numeric style (as in the project proposal). Overleaf runs biber
% automatically; locally use `latexmk -pdf main.tex`.
\usepackage[backend=biber,style=ieee,sorting=none,doi=true,url=false,isbn=false]{biblatex}
\addbibresource{references.bib}

\usepackage[hidelinks]{hyperref}
\usepackage[nameinlink,capitalise,noabbrev]{cleveref}
\crefname{proposition}{Proposition}{Propositions}

% =====================================================================
%  Notation (use these macros everywhere so the maths stays consistent)
% =====================================================================
\newcommand{\vocab}{\mathcal{V}}                 % vocabulary
\newcommand{\Yset}{\mathcal{Y}}                  % yes surface forms
\newcommand{\Nset}{\mathcal{N}}                  % no surface forms
\newcommand{\Dset}{\mathcal{D}}                  % rating digits 1..7
\newcommand{\pyes}{p_{\mathrm{yes}}}             % pooled conditional P(yes)
\newcommand{\capm}{\kappa}                       % captured mass
\newcommand{\ASIint}{\mathrm{ASI}_{\mathrm{int}}}% ASI_intrinsic
\newcommand{\PS}{\mathrm{PS}}                    % polarity separation
\newcommand{\RR}{\mathrm{RR}}                    % robustness rate
\newcommand{\CR}{\mathrm{CR}}                    % consistency rate
\newcommand{\dRR}{\Delta\mathrm{RR}}             % RR_vision - RR_text
\newcommand{\dm}{\delta_m}                       % extrinsic deviation
\newcommand{\rlow}{\bar r_{\text{low},q}}        % low-ASI human baseline
\newcommand{\rhigh}{\bar r_{\text{high},q}}      % high-ASI human baseline
\newcommand{\rexp}{\hat r}                       % expected rating
\newcommand{\rnorm}{\tilde r}                    % order-normalised rating
\newcommand{\ind}[1]{\mathbb{1}\!\left[#1\right]}
\newcommand{\E}{\mathbb{E}}
\DeclareMathOperator{\softmax}{softmax}
\DeclareMathOperator{\logit}{logit}
\DeclareMathOperator{\RMSNorm}{RMSNorm}

% =====================================================================
%  Planning helpers -- delete this block (and every plan box) at the end
% =====================================================================
% \showplantrue  -> plan boxes and figure placeholders are visible
% \showplanfalse -> they disappear (use this for the submitted PDF)
\newif\ifshowplan
\showplantrue

\usepackage[most]{tcolorbox}
\usepackage{environ}
\definecolor{planblue}{HTML}{1F4E79}
\definecolor{planfill}{HTML}{EEF4FA}
\definecolor{flagred}{HTML}{9C2A00}

\NewEnviron{plan}[1][Plan]{%
  \ifshowplan
  \begin{tcolorbox}[enhanced,breakable,colback=planfill,colframe=planblue,
      boxrule=0.6pt,arc=1mm,left=2mm,right=2mm,top=1mm,bottom=1mm,
      fonttitle=\bfseries\small,title={#1},fontupper=\small\singlespacing]
  \BODY
  \end{tcolorbox}
  \fi}

% Inline markers used inside plan boxes
\newcommand{\citeplan}[1]{\textbf{Cite:} \cite{#1}}
\newcommand{\pages}[1]{\textbf{Target length:} #1}
\newcommand{\decision}[1]{\textcolor{flagred}{\textbf{Decision needed:} #1}}
\newcommand{\newanalysis}[1]{\textcolor{flagred}{\textbf{New analysis:} #1}}
\newcommand{\source}[1]{\textit{Source: #1}}

% A known gap in the data, not a note to self. Shows as a red to-do while the
% plan boxes are on, and as an honest statement in the clean build, so a
% circulated PDF never shows "[fill]".
\newcommand{\gap}[1]{\ifshowplan\textcolor{flagred}{[\,#1\,?]}\else\emph{#1}\fi}

% Placeholder for a figure that has not been produced yet. It keeps the
% label, caption and List of Figures working, so cross-references compile.
%   \plannedfigure{label}{caption}{what goes in it / data / script}
\newcommand{\plannedfigure}[3]{%
  \begin{figure}[htbp]\centering
  \ifshowplan
    \fbox{\parbox{0.92\linewidth}{\small\singlespacing
      \textbf{Planned figure.} #3}}
  \else
    \fbox{\parbox{0.92\linewidth}{\centering\vspace{3cm}[figure]\vspace{3cm}}}
  \fi
  \caption{#2}\label{#1}
  \end{figure}}
